\begin{abstract}
When a model learns to generalise, does its behaviour become recoverable from one sparse circuit or from several? We follow five one-layer transformers through grokking on modular addition, searching at seven fixed checkpoints and evaluating every candidate over all 12,769 inputs. At 99\% agreement with the full model, greedy search retains 510--515 of 516 components in the delayed phase and 65--146 after stable generalisation. All fifteen stable positions support six or seven sparse masks; none of the twenty earlier positions yields a qualifying mask. A retrospective audit finds ten-member families among independently obtained non-monotone search endpoints at every stable position, without a diversity penalty during search. The original final families contain 27 circuits: none of 4,200 matched random masks reaches their fidelity. Their full neuron intersections are insufficient alone, yet removing them makes every member fail the fidelity criterion. Under calibrated Jensen--Shannon fidelity, stable endpoints retain 209--368 components, compared with 432--515 earlier, and four of fifteen stable positions support sparse families. The original prediction-preserving masks do not meet this stricter criterion. These results connect longitudinal compression to non-unique component realisations: generalisation accompanies a regime of several specific, compact selections that preserve decisions, while the tested searches recover larger masks when preserving output distributions.
\end{abstract}

\section{Introduction}

Grokking separates fitting the training data from learning a rule that generalises. A network can memorise an algorithmic dataset long before its test accuracy improves \citep{power2022grokking}. This interval makes grokking useful for studying how computation changes during learning. In modular-addition transformers, Fourier features can develop before the rise in test accuracy, followed by the removal of competing memorisation behaviour \citep{nanda2023progress}. In sparse-parity networks, prediction-preserving subnetworks become smaller as generalisation develops \citep{merrill2023two}. Compression is therefore already part of the empirical account of grokking.

Circuit discovery introduces a further question: how many sparse explanations can be recovered at each point in that trajectory? A procedure returning one mask cannot reveal whether it has found an isolated solution or one of several alternatives. Circuit non-uniqueness has been demonstrated in exhaustive small-network studies \citep{meloux2025everything}, and overlap-aware discovery can recover multiple faithful masks at trained endpoints \citep{chen2026rome}. Following such families during learning could distinguish a transition towards one selected explanation from a transition towards a regime with several recoverable explanations.

We study this distinction in five transformers trained on addition modulo 113. The task permits evaluation over the entire input domain, so a mask's measured fidelity does not depend on which evaluation examples happen to be sampled. We search the same basis of four attention heads and 512 MLP neurons at seven common training checkpoints. The original family criterion requires at least 0.990 agreement with the corresponding full model, no more than 258 retained components, and pairwise Jaccard overlap at most 0.50. Here a circuit is a behaviour-preserving component subset under zero ablation, and a recovered family is the set returned within a specified search budget. The study concerns the development of these recoverable explanations, not the number of distinct algorithms in the network.

The family result is sharp: every analysed stable position yields six or seven circuits, whereas every delayed and transitional position is empty at the primary criterion. The underlying change is also large. At 99\% prediction agreement, delayed first-search endpoints retain nearly all 516 components, while stable endpoints retain just 65--146. Generalisation is therefore associated with a change from almost full-model recovery to a regime in which several much smaller component selections reproduce the model's predictions. The transition recurs at different training steps in different models.

We establish this result through three connected measurements. First, a second, non-monotone optimiser recovers the size contrast and produces qualifying families without a diversity penalty. Second, matched random masks and shared-core interventions show that successful recovery depends on particular combinations of shared and alternative neurons. Third, distributional and margin fidelity reveal what the compact masks omit: much of the full model's confidence. The retained-size separation survives a strict distributional criterion, although sparse families then occur at fewer stable checkpoints.

Together, these measurements connect two observations usually studied separately: compression during learning and non-uniqueness of circuit explanations. Compression supplies the setting in which several sparse explanations can be recovered; multiplicity shows that the resulting explanation is not a unique component list. The size--fidelity relation describes how much of this freedom remains when the explanation must preserve richer output information. This longitudinal view makes circuit discovery a measurement of the developing model, rather than only a way to illustrate a trained endpoint.

\section{Related work}

\paragraph{Grokking and compression.}
\citet{nanda2023progress} traced Fourier-structured computation in modular-addition transformers, showing why useful internal structure need not first appear at the behavioural transition. \citet{merrill2023two} tracked prediction-preserving subnetwork size through sparse-parity grokking. \citet{manningcoe2026grokking} compared grokking with steady learning and followed feature encoding and compressibility through training. These studies establish compression across grokking as a direct antecedent. \citet{varma2023efficiency} explained grokking through competition between solutions of different efficiency, and \citet{demoss2025complexity} studied its rate--distortion dynamics. On modular addition itself, \citet{zhong2023clock} found different algorithms across models and parallel implementations within some models. Our addition is to ask when compact recovery admits several component realisations, whether that multiplicity survives removal of the diversity objective, and which information these alternatives preserve.

\paragraph{Circuit discovery and non-uniqueness.}
Automated circuit discovery searches for restricted computations that preserve selected behaviour \citep{conmy2023automated,bhaskar2024finding}. Faithfulness differs from component overlap \citep{hanna2024faith} and depends on the ablation methodology \citep{miller2024robust}. \citet{shi2024hypothesis} formalised tests of circuit faithfulness, localisation and minimality, providing a statistical precedent for evaluating recovered circuits against alternatives. \citet{meloux2025everything} demonstrated non-identifiability in small networks, while \citet{chen2026rome} deliberately searched for multiple faithful, low-overlap circuits. \citet{bayatmakou2026many} showed that structurally different circuits can share a sufficient core and remain causally interchangeable. Longitudinal circuit analysis also has precedent: \citet{tigges2024consistent} found algorithmic consistency across language-model training despite changes in individual components. We instead measure when several sparse selections become recoverable within each checkpoint, test their literal shared core, and compare preservation of decisions and confidence. Unlike trainable sparse-subnetwork studies \citep{frankle2019lottery}, our retained weights are fixed; no circuit is retrained.

\section{Experimental design}
\label{sec:methods}

\subsection{Models, checkpoints and components}

The task maps an ordered pair $(a,b)\in\{0,\ldots,112\}^2$ to $(a+b)\bmod113$. A fixed split places 3,830 pairs in training and 8,939 in testing. Each model is a one-layer causal transformer with residual width 128, four width-32 attention heads, a 512-neuron ReLU MLP, learned positional embeddings, and no normalisation or dropout. Five initialisations share the same data split. Training uses full-batch AdamW for 40,000 steps, with learning rate $10^{-3}$ and weight decay 1; details are in Supplement~S1.

All five models reach 99.9\% training accuracy by step 200. We classify a checkpoint as \emph{delayed} after memorisation and before the first saved checkpoint exceeding 10\% test accuracy; as \emph{transitional} thereafter until stable generalisation; and as \emph{stable} from the fifth checkpoint in the first sequence of five consecutive saved checkpoints reaching 99\% test accuracy. Checkpoints are saved every 50 steps. Stable onset occurs at steps 5,900, 9,050, 7,750, 10,400 and 6,850 for seeds 0--4.

Circuit analysis uses steps 200, 3,400, 7,450, 8,150, 8,500, 8,650 and 9,050. This common grid was taken from the pilot trajectory and frozen before the multi-seed family analysis. It contains nine delayed, eleven transitional and fifteen stable positions. Seed~3 becomes stable outside the grid, so the stable observations come from four models. Phase labels are assigned from the full-model training records, independently of circuit recovery.

The component universe $U$ contains four whole attention heads and 512 MLP neurons. For a retained mask $C\subseteq U$, excluded head outputs are zeroed before entering the residual stream and excluded neuron activations are zeroed after ReLU, before the MLP output projection. Embeddings, unembeddings, biases and all other operations remain active. We measure size as $|C|$, with the original sparsity convention $|C|\leq258$.

Masking changes the computation without fitting a replacement model. The retained weights are exactly those of the checkpoint, and the same mask is applied to every input and token position. A successful mask therefore demonstrates that the checkpoint already supports the measured behaviour through that component selection. The all-retained mask is checked against the unmodified model before analysis; hook-restoration and parameter-immutability checks ensure that evaluating one candidate does not change later candidates.

\subsection{Fidelity and family discovery}

For full checkpoint $M$, masked checkpoint $M_C$, and the complete input domain $X$, top-one fidelity is
\begin{equation}
 F(C;M)=\frac{1}{|X|}\sum_{x\in X}
 \mathbf{1}\{\arg\max M_C(x)=\arg\max M(x)\}.
 \label{eq:fidelity}
\end{equation}
This is agreement with the checkpoint, not accuracy against arithmetic labels. Candidate masks are evaluated on all $|X|=\N$ inputs. A family contains masks satisfying $F\geq\tau$, $|C|\leq258$, and
\begin{equation}
 J(C_i,C_j)=\frac{|C_i\cap C_j|}{|C_i\cup C_j|}\leq\delta
 \quad\text{for all }i\ne j.
\end{equation}
The primary values are $\tau=0.990$ and $\delta=0.50$. The complete sensitivity grid crosses six fidelities, $0.800,0.850,0.900,0.950,0.975,0.990$, with three overlap ceilings, $0.25,0.50,0.75$, at all 35 positions: 630 family-search settings.

The first search starts with every component retained. A Taylor estimate of deletion damage ranks candidates; exact evaluation determines acceptance. Within the first candidate batch containing an admissible deletion, the highest-fidelity candidate is selected. Subsequent family members are sought with a component-reuse penalty and must pass the pairwise overlap constraint. Each requested member has an allowance of 10,000 exact evaluations, each family cell has 50,000, and the target family size is ten. The first member follows one deterministic path; later members may use up to five seeded restarts. Supplement~S2 gives the algorithm and stopping rules.

We also recover the terminal size of the first path before applying the 258-component filter. This yields 210 distinct seed--checkpoint--fidelity endpoints, because the overlap ceiling cannot affect the first member. These endpoints make the size trajectory observable even where the reported family is empty.

A fidelity of 0.990 allows at most 127 prediction disagreements across all ordered input pairs. Exact evaluation is used both to accept deletions and to report final masks. Retained size measures the compression achieved by one path; family size measures the recovery of alternatives. Reporting both distinguishes an empty family whose first endpoint is nearly full-sized from one narrowly missing the sparsity cutoff.

\subsection{Search and fidelity robustness}

The independent search proposes uniformly sampled deletions, additions and swaps with probabilities $0.5,0.1,0.4$. It starts from the full mask and rejects any proposal violating the active fidelity constraint. Feasible deletions and swaps are accepted; feasible additions use a cooled Metropolis rule with retained count as the energy. This permits the search to replace components and revisit larger masks. Each of ten streams per position receives up to 20,000 unique stochastic evaluations or 40,000 proposals, followed by a separate canonical-order deletion cleanup. A stable positive-control gate preceded the full top-one analysis.

The two searches differ in how they explore component choices. Greedy deletion uses a gradient-derived ranking and commits to every removal. Annealing uses no Taylor ranking: a swap can exchange one retained component for an excluded one without changing size, and an addition can reopen a deletion route unavailable to the current mask. The final cleanup then applies a common, readily checked stopping condition: every remaining single deletion fails the active fidelity constraint. Ten separately seeded streams expose variation across these search paths at each fixed checkpoint.

For non-argmax fidelity, let $p_x=\operatorname{softmax}(M(x))$, $q_x=\operatorname{softmax}(M_C(x))$, and $r_x=(p_x+q_x)/2$. We use mean Jensen--Shannon divergence
\begin{equation}
 \JS(C;M)=\frac{1}{2|X|}\sum_x
 \left[\operatorname{KL}(p_x\|r_x)+\operatorname{KL}(q_x\|r_x)\right],
 \label{eq:js}
\end{equation}
with natural logarithms and temperature~1. We separately measure the mean absolute change in the full model's top-class logit margin, normalised by its mean absolute reference margin (Supplement~S3). Both criteria use complete-domain evaluation.

These criteria ask for different kinds of preservation. Top-one agreement keeps the winning class but permits its probability to change. Jensen--Shannon divergence compares the full probability vectors, including non-winning classes. Margin error instead keeps track of the gap between the reference class logit and its strongest competitor, even when softmax probabilities are already close to one. For example, two outputs can share their winning class while assigning it very different confidence; sufficiently large logits can also change substantially with little change in softmax probability. Measuring all three distinguishes compact recovery of decisions from compact recovery of distributions or logit separation. Their numerical tolerances are not interchangeable units of fidelity.

The deterministic first-circuit search is repeated at six tolerances per criterion, giving 420 runs. Its Taylor ranking is held fixed; exact admissibility uses the new criterion. For annealing, a prespecified stable-only calibration selects the strictest tolerance at which at least four of five streams recover at most 258 components, all complete cleanup, and none is censored. Jensen--Shannon selects $0.0005$; no margin tolerance passes. We therefore run the full 350-stream annealing analysis for Jensen--Shannon only, while retaining the entire deterministic grid for both metrics.

\subsection{Matched nulls and evidence structure}

At step 9,050, the four stable primary families supply 27 circuits and 78 within-family pairs. For each pair we compare observed Jaccard overlap with independent uniformly drawn subsets preserving either each mask's total size or its separate head and neuron counts. Exact hypergeometric calculations provide the null expectation, interval and tails; deterministic Monte Carlo draws check those calculations. For behavioural nulls, we evaluate 100 masks per unique seed/composition profile under each matching rule: 21 profiles and 4,200 full-domain evaluations.

The independently trained model is the unit of replication. Checkpoints, masks, pairs and optimiser streams are nested observations. We report their distributions descriptively and retain the original seed-level paired endpoint in Section~\ref{sec:families}. The family analysis was frozen before these follow-ups. Structural and behavioural nulls and first-path reconstruction (E1--E3) were designed after the original results; search and non-argmax follow-ups (E4--E5) had prospective rules before their respective outcomes. Supplement~S4 records the sequence and calibration decisions.

An additional exploratory audit tests family compatibility among the ten saved annealing masks at each stable position, without rerunning or changing their search objectives. Exhaustive enumeration gives the largest qualifying subset of each ten-mask pool. A separate evaluation-only follow-up uses the original 27 final family members. With all four heads retained, we evaluate their neuron intersection, their union, and each member after removing the intersection. We also score the same masks under JS at common temperatures 1, 2 and 4 and under reference-margin error. These checks have a written protocol, fixed coverage and separate provenance (Supplement~S7).

\section{Results}

\subsection{The change in retained size is much larger than the sparsity cutoff}
\label{sec:trajectory}

At fidelity 0.990, the first greedy path ends with 510--515 components during the delayed phase, 297--515 during transition, and 65--146 after stable generalisation. The corresponding medians are 514, 514 and 82. Every endpoint is locally single-deletion minimal and uncensored. Thus the primary family boundary does not arise from endpoints narrowly crossing 258: delayed masks remain almost full-sized, while stable masks retain less than 29\% of the component basis (Figure~\ref{fig:trajectory}).

\begin{figure}[tb]
\centering
\includegraphics[width=\linewidth]{trajectory.pdf}
\caption{\textbf{Large reductions in recovered size accompany stable generalisation.} (a) Unsmoothed test accuracy for five models; marks below the axis show the seven circuit-analysis steps. (b) Every first greedy endpoint at fidelity 0.990. Each cell gives the retained count out of 516 and its independently assigned phase: D, delayed; T, transitional; S, stable. Darker cells retain fewer components. Columns are the common sampled checkpoints, not equally spaced training times. All fifteen stable endpoints lie well below the 258-component cutoff; all twenty earlier endpoints lie above it.}
\label{fig:trajectory}
\end{figure}

The separation persists across the six top-one fidelity levels. Even at 0.800, delayed endpoints retain 483--496 components and stable endpoints 39--64. No delayed endpoint meets the size cutoff at any threshold; every stable endpoint does. Transition is graded: 4/11 positions meet it at 0.800, 2/11 at 0.850 and 0.900, 1/11 at 0.950 and 0.975, and none at 0.990. These are separate threshold searches, not a fitted continuous curve.

The common checkpoint grid also separates elapsed training from behavioural phase. At step 7,450, the stable seeds 0 and 4 already admit compact masks, while seeds 1, 2 and 3 are still transitional. Seed~2 subsequently joins the compact regime, and seed~1 does so only at the final analysed step. Seed~3 remains transitional throughout this grid. Thus a late checkpoint is not sufficient for compact recovery, and the result is not driven by assigning all models the same transition time. The phase labels come from the original accuracy trajectories, not from where retained size happens to fall.

The independent search reaches a similar size separation (Figure~\ref{fig:robustness}). Across 350 top-one annealing runs, delayed endpoints retain 508--515 components, transitional endpoints 264--515, and stable endpoints 58--136. All 150 stable runs meet the original sparsity criterion and all 200 earlier runs fail it. Addition and swap moves therefore do not remove the contrast found by greedy deletion. Budget censoring remains concentrated in the earlier phases; Table~S3 reports it separately from the local cleanup checks.

\begin{figure}[tb]
\centering
\includegraphics[width=\linewidth]{search_robustness.pdf}
\caption{\textbf{Changing the search preserves the retained-size contrast.} Each panel pairs greedy and annealing endpoints at the same 35 seed--checkpoint positions, using the same fidelity criterion on both axes. Points give the median of ten annealing streams; vertical bars span their full range, not a confidence interval. The diagonal denotes equal recovered size; points below it indicate smaller annealing masks. Dashed lines mark 258 components. Stable positions occupy the low-size region under both criteria, but many distribution-preserving endpoints exceed 258. Coincident positions can overlap. All endpoints complete local deletion cleanup; budget and plateau counts are reported in Table~S3.}
\label{fig:robustness}
\end{figure}

\subsection{Stable models admit multiple behaviourally special masks}
\label{sec:families}

Under the primary family criterion, none of the nine delayed or eleven transitional positions yields an eligible family. All fifteen stable positions yield six or seven circuits (Table~\ref{tab:families}). At step 7,450, for example, seeds 0 and 4 are stable and each yields seven circuits, while seeds 1, 2 and 3 are transitional and yield none. Seed~2 enters the recovered-family regime by step 8,150 and seed~1 by step 9,050. The association therefore recurs at different training steps.

\begin{table}[tb]
\centering\small
\caption{\textbf{Primary family size on the common checkpoint grid.} D, T and S denote delayed, transitional and stable phases. Entries give phase and recovered member count.}
\label{tab:families}
\begin{tabular}{crrrrrrr}
\toprule
Seed & 200 & 3,400 & 7,450 & 8,150 & 8,500 & 8,650 & 9,050\\
\midrule
0 & D:0 & T:0 & S:7 & S:7 & S:7 & S:6 & S:6\\
1 & D:0 & D:0 & T:0 & T:0 & T:0 & T:0 & S:7\\
2 & D:0 & D:0 & T:0 & S:7 & S:7 & S:7 & S:7\\
3 & D:0 & D:0 & T:0 & T:0 & T:0 & T:0 & T:0\\
4 & D:0 & D:0 & S:7 & S:7 & S:6 & S:6 & S:7\\
\bottomrule
\end{tabular}
\end{table}

Across all 630 family settings, delayed positions are nonempty in 0/162 cells, transitional positions in 30/198, and stable positions in 270/270. Tightening the overlap ceiling to 0.25 reduces the number of alternatives but still yields three to seven circuits in each final stable family at fidelity 0.990. The original paired endpoint, the change in primary family size from step 200 to 9,050, is $+6,+7,+7,0,+7$ across seeds. Its two-sided exact sign-test probability is 0.125 after excluding the zero difference. The repeated checkpoint surface describes where the change occurs; it does not increase the number of independent models.

At the final checkpoint, the four stable models yield 27 masks retaining 66--180 components, with fidelity 0.9901--0.9918. All retain the four attention heads. Their 78 within-family Jaccard values range from 0.0657 to 0.3092. Size matching changes the meaning of these small raw overlaps: the pooled observed median is 0.1581, whereas the median pair-specific expectation is 0.0879 under total-size matching and 0.1085 when head and neuron counts are preserved. The families are more overlapping, not more disjoint, than those random-subset baselines (Supplement, Figure~S1).

Their behaviour is markedly different from random subsets. None of the 4,200 matched random masks reaches fidelity 0.990. Under total-size matching, median fidelity is 0.0261 and the maximum is 0.4120. Preserving the four heads and the exact neuron count produces a stronger null, with median 0.2830 and maximum 0.6361, still far below every recovered circuit (Figure~\ref{fig:nulls}). Low overlap alone is not the evidence that makes these masks interesting. It is the recovery of several different component selections that each preserve behaviour where typical selections of the same size and composition do not.

The composition-matched comparison is important because all 27 recovered masks retain every attention head. A total-size null often removes heads, so its failure could largely reflect the loss of attention rather than a consequential choice among MLP neurons. The stronger null leaves attention intact and changes only which neurons survive, retaining exactly the observed number. Its substantial improvement over total-size matching confirms that this distinction matters. Yet even its best draw disagrees with the full model on more than a third of the domain, compared with less than one percent for each recovered circuit. Correct head retention and a generous neuron count alone do not reproduce the recovered behaviour.

Multiplicity also appears without a diversity objective. At each of the fifteen stable positions, all ten saved top-one annealing masks meet the original size, fidelity and pairwise-overlap criteria. The independent streams therefore supply a ten-member family at every stable position, rather than merely ten repeated recoveries of the same mask. This retrospective check uses the original overlap cutoff; at the stricter 0.25 cutoff, the largest compatible subset ranges from one to ten. The original six or seven members were consequently not an upper bound on available alternatives.

The shared-core evaluation tests how the literal common neurons contribute to each family. Their all-member intersections contain only 3, 8, 4 and 1 MLP neurons in seeds 0, 1, 2 and 4. With all heads retained, these intersections reach fidelity 0.032--0.046; removing them from each original member lowers fidelity to 0.692--0.987, below 0.990 in all 27 cases (Table~\ref{tab:core}). The shared set therefore contributes collectively to every tested member, while the full intersection alone does not reproduce its predictions. Fidelity is not monotone in mask inclusion: even the union of all members fails the threshold in three models. Combining faithful selections can disturb their successful behaviour.

\begin{table}[tb]\centering\small
\caption{\textbf{Shared neurons matter, but do not suffice.} All interventions retain the four attention heads. Core: neurons in every original family member. Entries report full-domain prediction agreement. Removing the core fails the 0.990 criterion for all 27 members; taking the union of members also fails in three models.}
\label{tab:core}
\begin{tabular}{crrrrr}\toprule
Seed & Members & Core neurons & Core only & Core removed & Union \\\midrule
0 & 6 & 3 & 0.0458 & 0.7723--0.9511 & 0.9829 \\
1 & 7 & 8 & 0.0341 & 0.6918--0.7622 & 0.9995 \\
2 & 7 & 4 & 0.0408 & 0.7710--0.8502 & 0.8599 \\
4 & 7 & 1 & 0.0324 & 0.9290--0.9874 & 0.8541 \\
\bottomrule\end{tabular}\end{table}


\begin{figure}[tb]
\centering
\includegraphics[width=\linewidth]{matched_nulls.pdf}
\caption{\textbf{Recovered circuits preserve predictions where matched random masks fail.} At step 9,050, all 27 recovered circuits reach 0.990 agreement (diamonds, nearly coincident within each seed). Neither random-mask null reaches that line in any of 4,200 evaluations. Boxes show the middle 50\% of random fidelities, centre lines the median, and whiskers the full observed range. Random draws are pooled within each model with 100 draws per distinct matching profile and null; profiles shared by several recovered circuits are counted once. Blue boxes preserve all four heads and the exact neuron count; grey boxes preserve only total size. These are random-mask distributions, not confidence intervals across trained models.}
\label{fig:nulls}
\end{figure}

The original operand-subset probe provides a descriptive check on the recovered families. Splitting the test inputs by low and high values of each operand gives four fidelity measurements per circuit. All fifteen stable families form one group at the original 0.05 profile-distance tolerance; their largest within-family coordinate difference is below 0.01. Structural variation therefore does not produce a large separation under this coarse behavioural partition. In a one-seed Fourier pilot, the full model and one circuit at each fidelity level also retain the expected addition signature; the primary 146-component circuit places 0.7316 of non-DC logit power on the addition manifold. These secondary results are consistent with alternative component selections preserving the same task-level behaviour. Their full coverage and definitions are given in Supplement~S6.

\subsection{Distribution-preserving searches recover larger circuits}
\label{sec:nonargmax}

The non-argmax analyses show that substantial compression and crossing the half-model threshold are different results. At the calibrated Jensen--Shannon tolerance $0.0005$, the independent search retains 209--368 components at stable positions, versus 432--515 at transitional positions and 509--515 when delayed. Every observed stable endpoint is smaller than every observed earlier endpoint. The stable median is 276.5, however, and only 40/150 stable streams satisfy $|C|\leq258$ (Figure~\ref{fig:robustness}, right). Those 40 runs come from four positions, each with ten successes: seed~0 at step 9,050, seed~2 at step 9,050, and seed~4 at steps 7,450 and 9,050.

These are multi-mask outcomes under distributional fidelity as well. At Jaccard 0.50, the saved pools contain ten-member families at the three final positions and a four-member family at seed~4, step 7,450. At 0.25, none contains a qualifying pair. The fixed overlap ceilings define compatibility, not unusual disjointness relative to chance; larger masks also have higher size-matched expected overlap. Under the calibrated criterion, the observed multiplicity occurs at fewer positions and with greater overlap. One of these four positions, seed~0 at step 9,050, supplied the separate calibration runs.

All stable Jensen--Shannon runs are uncensored. All 90 delayed runs and 48/110 transitional runs are budget-censored, although every run finishes local deletion cleanup. The endpoint ranges establish separation in what these searches recovered, while the search records show where further exploration remains most consequential.

The full deterministic tolerance grids make the metric dependence explicit (Supplement~S3). At Jensen--Shannon tolerance 0.0005, the median sizes are 515, 514 and 281 in delayed, transitional and stable positions; 5/15 stable positions meet the size cutoff. At tolerance 0.005, all fifteen stable positions do, with median size 191, while no earlier position does. No Jensen--Shannon tolerance yields a sparse delayed endpoint.

Figure~\ref{fig:fidelity} places this result within the complete tolerance grids. Each panel moves from strict to permissive preservation. The stable median lies below the earlier-phase medians throughout all three grids, while its absolute retained size changes substantially with the metric and tolerance. The transition is therefore visible as a size difference even at settings where no phase meets the binary sparsity rule. At the other extreme, loosening the criterion allows additional deletion without making the delayed models comparably compact. The full grids retain both effects rather than selecting a single tolerance that produces the cleanest family boundary.

\begin{figure}[tb]
\centering
\includegraphics[width=\linewidth]{fidelity_tradeoff.pdf}
\caption{\textbf{The fidelity requirement determines how much compression is recoverable.} Complete deterministic first-search grids for top-one agreement, Jensen--Shannon divergence and normalised reference-margin error. Each line shows the median across positions in a phase; shading spans their minimum to maximum, not estimation uncertainty. There are nine delayed, eleven transitional and fifteen stable positions at every tolerance. Each plotted point comes from a separate search; connecting lines guide the eye. Lower top-one thresholds and higher discrepancy tolerances are more permissive. One margin endpoint is budget-censored and retained as observed (Table~S4). Dashed lines mark 258 components.}
\label{fig:fidelity}
\end{figure}

Margin preservation is more demanding still. At normalised margin error 0.20, median retained sizes are 507, 480 and 349 in the three phases, with no sparse endpoint. At 0.40 they are 493, 426 and 246, and 11/15 stable positions are sparse. In the stable-only annealing calibration, even this loosest margin tolerance yields only 3/5 sparse masks, below the required four successes. The full margin annealing analysis is therefore not run. The margin grid remains evidence for greater post-generalisation compressibility, but not for uniformly sparse recovery under stringent preservation of logit margins.

Evaluating the original 27 masks across metrics makes the information loss concrete. Although every mask preserves at least 99\% of predictions, their JS divergences are 0.0252--0.2130 and their normalised margin errors 0.653--0.889. None meets either the calibrated JS tolerance or even the loosest margin tolerance. Seven masks, all from seed~1, do meet the loosest JS tolerance of 0.05. At common temperatures 2 and 4, JS rises to 0.116--0.463 and 0.218--0.487: softening the distributions exposes differences partly hidden by confident decisions. The compact decision masks and the larger distribution-preserving masks are not interchangeable evidence.

Confidence also covaries with longitudinal recovery. Across the fifteen stable positions, mean reference margin and median JS-annealing size have descriptive Spearman correlation $-0.949$; within seeds 0, 2 and 4 the correlations are $-1.00,-1.00,-0.90$. Seed~1 supplies only one stable position. These nested trajectories do not identify a confidence-adjusted effect, but they make confidence a supported explanation for part of the JS compression, rather than an untested possibility.

\section{Discussion}

The main finding is a change in the recoverability of model behaviour across grokking. Before stable generalisation, the tested searches usually retain nearly the full component basis to preserve the checkpoint's predictions. Afterwards, extensive deletion is possible and the original family search finds several alternative masks. This combination matters: compression establishes that a small component selection can suffice, while family recovery establishes that the sufficient selection is not unique. The matched random-mask controls show that the alternatives are specific enough to preserve behaviour that arbitrary matched selections lose.

The strongest longitudinal evidence is the magnitude and consistency of the retained-size difference. Stable and earlier endpoints remain separated after replacing gradient-ranked deletion with non-monotone search, and again when the new search preserves output distributions. These comparisons probe different dependencies of the original result: deletion history in one case, the information preserved by fidelity in the other. They support a compression transition that is broader than the original choice of a 258-component boundary. The complete checkpoint map also shows this transition recurring at different times, rather than following elapsed training alone.

The fidelity results give this transition more structure. Stable models are particularly compressible when the target is their discrete decision. At the tested tolerances, searches preserving output distributions return larger masks, and margin-preserving endpoints can be larger still. These are complementary measurements of the same checkpoint, not competing verdicts about whether it contains a circuit. A compact mask suitable for explaining predictions may omit information needed to explain confidence. Reporting the retained-size relation makes that difference explicit and prevents a thresholded family diagram from hiding it.

One plausible account is that learning the arithmetic rule makes behaviour easier to reproduce through several component selections. Before generalisation, the reference includes irregular predictions away from the training set; afterwards, most predictions follow one rule. At the same time, training can change how that rule is represented and how readily search can isolate it. The present evidence measures the resulting recovery problem. It is compatible with a useful Fourier-structured computation developing earlier, as in prior grokking studies, and becoming more separable from the rest of the model as generalisation develops. Distinguishing these contributions is a mechanistic question beyond the observed compression trajectory.

The family interventions give this non-uniqueness a concrete form. Removing the all-member intersection disrupts every tested selection, yet retaining that intersection alone is insufficient. Nor are faithful masks closed under taking unions. These observations identify contributions that depend on the surrounding selection. Ablation-induced compensation and counterbalancing are known in language models \citep{mcgrath2023hydra}; our union test demonstrates non-monotone fidelity, without identifying those mechanisms here. Combined with the random-mask failures and the independently recovered families, the interventions show that mask identity matters beyond retained size. Following these alternatives through training reveals when the tested searches recover several compact realisations of the model's decisions.

\section{Limitations}

The full checkpoint's behaviour is the reference at every training step. Before generalisation, fidelity requires preserving its errors as well as its correct outputs; afterwards, it largely requires preserving the arithmetic rule. Thus the observed size change could reflect a simpler prediction function, a computation better aligned with the component basis, an easier search landscape, or a combination. It does not locate the first appearance of a generalising mechanism or establish a causal effect of grokking. The one available random-label control yields no primary family at seven checkpoints, but changes task structure as well as generalisation. A planned structured no-generalisation control did not meet its selection rule.

Both searches operate on the same mixed head/neuron basis and use zero ablation. A head and a neuron have different parameter and computational costs, so retained count is not a parameter- or FLOP-compression ratio. Annealing can add and swap components but never accepts an infeasible mask; it cannot traverse arbitrary disconnected feasible regions. Nor does finishing the allocated search establish convergence: only 43/150 stable top-one streams and 21/150 stable Jensen--Shannon streams meet the recorded stochastic plateau criterion, despite all completing local cleanup. Locally minimal endpoints and empty families are therefore not global optima or nonexistence results. The retrospective compatibility audit establishes multiplicity within the independently obtained pools, not saturation of the family space.

The confidence association is not resolved by evaluating fixed masks at other temperatures: that would require a controlled longitudinal search. Because fidelity is non-monotone, failure of the full intersection does not exclude a sufficient proper subset; neither those subsets nor approximately shared cores were exhaustively tested. The interventions also do not test causal interchangeability. Likewise, the four operand subsets were held out from training but included in full-domain circuit search. Their profiles are descriptive subdivisions of the discovery domain, not independent functional validation. Neither these tests nor the single-seed Fourier pilot establish distinct algorithms.

Finally, the study uses one task, one architecture and five initialisations on one split, with four models contributing stable checkpoints. The grid samples late training unevenly and omits seed~3's stable onset. More broadly, the robustness studies follow an observed original result, although their execution rules were frozen before the new search outcomes. The data support a consistent pattern in this setting, not a population-level law about grokking.

\section{Conclusion}

Sparse circuit recovery changes substantially across the analysed grokking trajectories. Under top-one fidelity, stable checkpoints support several compact, behaviourally special masks; earlier checkpoints do not yield qualifying masks. Changing both search and fidelity preserves a large separation in recovered size but not a universal half-model sparsity boundary. The useful object is therefore not a family count in isolation, but the family together with the size--fidelity relation under which its members become recoverable.

\section*{Reproducibility}
The supplement specifies the search algorithm, budgets, metrics, calibration decisions and complete checkpoint summaries. The original experiment is archived at \url{https://doi.org/10.5281/zenodo.21917638}. The \href{https://github.com/AK-314/circuit-families/releases/tag/phase1-tmlr-review-followup-2026-09-13}{versioned follow-up release} supplies E1--E5 code and tables, saved search endpoints, the exploratory evaluator and its per-input results, and the manuscript sources. File manifests preserve the original freeze and identify every new input and output. The included instructions separate figure reconstruction, saved-mask evaluation and full search reproduction.

\section*{Use of AI}
Generative AI assisted with coding, analysis, manuscript drafting, editing and formatting. The author is responsible for checking the analyses, references and final scientific claims.

\section*{Acknowledgements}
I thank Eleftheria Fassman of Cornell University for advice and guidance throughout this project.

\bibliographystyle{tmlr}
\bibliography{references}
